\item Al pasar por la materia, la intensidad $I$ de un haz de fotones (medido en $\text{W/m}^2$) disminuye exponencialmente según la ley de atenuación:
$$ I = I_0 e^{-\mu x} $$
donde $\mu$ es el coeficiente de absorción lineal y $x$ es el espesor del material.
\begin{enumerate}
    \item Demuestre que si dos haces con longitudes de onda $\lambda_1$ y $\lambda_2$ e intensidades incidentes idénticas pasan por una misma placa, la razón de sus intensidades emergentes es:
    $$ \frac{I_2}{I_1} = e^{-(\mu_2 - \mu_1)x} $$
    \item Calcule esta razón para un haz emergente de una placa de aluminio de $1.00\text{ mm}$ de espesor si el haz incidente contiene intensidades iguales de rayos X de $50\text{ pm}$ ($\mu_1 = 5.4\text{ cm}^{-1}$) y $100\text{ pm}$ ($\mu_2 = 41.0\text{ cm}^{-1}$).
    \item Repita el cálculo anterior utilizando una placa de aluminio de $10.0\text{ mm}$ de espesor.
\end{enumerate}